\chapter{出版说明}

当前，中学教学改革已经深入到课程设置和教材改革领域。
我校数学教材的改革，以发展学生的数学思维为目标，以不改变现行教学大纲规定的教学闪容为前提，试图通过对知识结构及其展开方式的统盘考虑，实现整体优化。经多年反复探索、实验，编成了这套尝试融教材与教法、学法于一体的《北京四中高中教学讲义》。

这套讲义的产生可以上溯到 1982 年。从那时起，为了发展学生智能，提高数学素养，我校部分同志就开始对高中数学教学进行以教材改革为龙头，以学法教育为重点的“整体优化实验研究”。正是在这项研究的基础上，逐步形成了这套讲义编写的特色和风格。这就是：
\begin{enumerate}
    \item 为形成学生良好的认知结构，讲义的知识结构力求脉络分明，使学生能从整体上理解教材。
    \item 为了提高学生的数学素养，本讲义把数学思想的阐述放到了重要位置。数学思想既包含对数学知识点（概念、定理、公式、法则和方法）的本质认识，也包含对问题解决的数学基本观点。它是数学中的精华，对形成和发展学生的数学能力具有特别重要的意义。为此，讲义注重展现思维过程（概念、法则被概括的过程，教学关系被抽象的过程，解题思路探索形成的过程）。在过程中认识知识点的本质，在过程中总结思维规律，在过程中揭示数学思想的指导作用。力图使学生能深刻领悟教材。
    \item “再创造，再发现”在数学学习中对培养创造维能力至关重要，为引导学生积极参与“发现”，讲义在设计上做了某些尝试。
    \item 例题和习题的选配，力求典型、适量、成龙配套。习题分为 A 组（基本题）、 B 组（提高题）和 C 组（研究题）。教师可根据学生不同的学习水平适当选用。
    \item 教材是学生学习的依据。应有利于培养自学能力，本书注重启迪学法，并在书末附有全部习题的答案或提示，以供学习时参考。
 \end{enumerate}   

这套讲义在研究、试教和成书的过程中，始终得到了北京市和西城区教育部门有关领导的关怀和帮助，得到了北京师范大学数学系钟善基教授、曹才翰教授的热情指导，清华附中的瞿宁远老师也积极参与了我们的实验研究，并对这套教材做出了贡献，在此一并致以诚挚的谢意。

在编写过程中，北京四中数学组的教师们积极参加研讨，对他们的热情支持表示感谢。

这套讲义包括六册：高中代数第一、二、三册，三角、立体几何、解析几何各一册。

编写适应素质教育的教材，对我们来说是个尝试。由于水平所限，书中不当之处在所难免，诚恳希望专家、同行和同学们提出宝贵意见。

\begin{flushright}
    北京四中教学处\\
    1996年1月
\end{flushright}

\chapter{绪论}
17世纪上半叶，数学正处在由“常量数学”向“变量数学”的发展时期. 在这一巨大变革时期，数学科学逐渐出现了两个崭新的分支——解析几何学和微积分学. 解析几何学的出现，从本质上改变了数学的面貌. 正如著名数学家拉格朗日（Lagrange）指出的：“只要代数同几何分道扬镳，它们的进展就缓慢，它们的应用就狭窄. 但是当这两门科学结合成伴侣时，它们就互相吸取新鲜的活力，从那以后，就以快速的步伐走向完善. ”

解析几何学的产生不是偶然的，是有着深刻的历史背景. 16—17世纪，欧洲大陆由封建社会向资本主义社会过渡，生产力得到极大的解放，促进了科学技术的飞速发展. 这不仅为解析几何学的产生和发展提出了需求，而且为它的产生和发展提供了条件. 同时，几何学和代数学（主要是方程理论）的发展，促使人们对数学方法论的研究产生了极大的兴趣. 人们不再满足欧几里得几何学的公理体系的完美和方法的巧妙，正像数学家笛卡尔（Descartes）写道：“我决心放弃那个仅仅是抽象的几何. 这就是说，不再去考虑那些仅仅是用来练习思想的问题. 我这样做，是为了研究另一种几何，即目的在于解释自然现象的几何. ”

解析几何学的建立，是与数学家费尔马（Fermat）和笛
卡尔的卓有成效的研究分不开的. 特别是笛卡尔提出两个基本观念——用坐标表示点的观念和用方程表示曲线的观念，为解析几何学的发展打下了坚实的基础. 

解析几何学的目标是借助于代数方法来研究几何的问题. 因此，下述问题就成为解析几何学的\textbf{两个基本问题}：一是，已知曲线求其方程，再通过方程研究曲线的性质；二是，已知方程画出其曲线，再用曲线研究方程的性质. 

鉴于解析几何学本身的特点，在学习时应注意以下两点：第一，“数”“形”结合，即看到“形”时联想“数”，研究“数”时想到“形”；第二，变量的思想，即用运动、变化的观点研究问题. 

解析几何学分为平面解析几何学与空间解析几何学. 高中阶段仅限于学习平面解析几何学. 为了有利于对解析几何学本质的理解，本书第一章集中介绍了解析几何学的基础. 第二、三章专题研究平面一次、二次曲线. 第四、五章学习坐标变换和极坐标. 第六章深入研究解析几何学的一些基础问题. 由于参数方程的思想在理论上和实践中有一定的重要性，因此，在第三章中我们就对有关概念做了介绍，同时在研究各种具体曲线时，都适时引入了它们的参数方程，并初步介绍了参数方程的应用. 

有“*”号的部分为选学内容. 






